\begin{mdframed}[style=mdpurplebox,frametitle={IMO 2025 Problem 1}]
A line in the plane is called \emph{sunny}
if it is not parallel to any of the $x$-axis, the $y$-axis, or the line $x+y=0$.

Let $n \ge 3$ be a given integer.
Determine all nonnegative integers $k$ such that there exist $n$ distinct lines
in the plane satisfying both of the following:
\begin{itemize}
\ii for all positive integers $a$ and $b$ with $a+b\le n+1$,
  the point $(a,b)$ lies on at least one of the lines; and
\ii exactly $k$ of the $n$ lines are sunny.
\end{itemize}
\end{mdframed}
\solutionheading{Solution}
\begin{remark*}
  7/7.
\end{remark*}
\textbf{Definition.} For an integer \(n\ge 3\) let\par
\[
P_n = \{(a,b)\in\mathbb{N}^2 \mid a\ge 1,\; b\ge 1,\; a+b\le n+1\}.
\]
We are asked for which nonnegative integers \(k\) there exist \(n\) distinct lines in the plane such that\par
\compactbullet{every point of \(P_n\) lies on at least one of the lines, and}
\compactbullet{exactly \(k\) of the lines are \emph{sunny} (i.e. not parallel to the \(x\)-axis, the \(y\)-axis, or the line \(x+y=0\)).}
We will prove that for every \(n\ge 3\) the only possible values of \(k\) are \(0,1,3\).\par
\smallskip
\solutionheading{1. Constructions for \(k = 0,1,3\)}
We give explicit families of \(n\) distinct lines that satisfy the two conditions.\par
\solutionheading{\(k = 0\)}
Take the \(n\) horizontal lines\par
\[
y = 1,\; y = 2,\; \dots,\; y = n.
\]
Each point \((a,b)\in P_n\) lies on the line \(y = b\). No line is sunny (all are parallel to the \(x\)-axis). Hence \(k = 0\) is attainable.\par
\solutionheading{\(k = 1\)}
Take the \(n-1\) horizontal lines\par
\[
y = 1,\; y = 2,\; \dots,\; y = n-1.
\]
These cover all points with \(b \le n-1\). The only point not covered is \((1,n)\) (since \(a+n\le n+1\) forces \(a = 1\)). Add one extra line, for instance\par
\[
L:\; y = x + (n-1).
\]
\(L\) has slope \(1\), so it is sunny. It passes through \((1,n)\). The total number of lines is \((n-1)+1 = n\), and exactly one line is sunny. Thus \(k = 1\) is attainable.\par
\solutionheading{\(k = 3\)}
\emph{If \(n = 3\):}Use the three sunny lines\par
\[
\begin{aligned}
L_1 &: y = x,\\
L_2 &: y = -2x + 5,\\
L_3 &: y = -\tfrac12 x + \tfrac52.
\end{aligned}
\]
Check: \(L_1\) contains \((1,1)\) and \((2,2)\); \(L_2\) contains \((1,3)\) and \((2,1)\); \(L_3\) contains \((1,2)\) and \((3,1)\). The six points of \(P_3\) are all covered. All three lines are sunny, so \(k = 3\).\par
\emph{If \(n > 3\):}First take the \(n-3\) horizontal lines\par
\[
y = 1,\; y = 2,\; \dots,\; y = n-3.
\]
These cover all points with \(b \le n-3\). The remaining points are those with \(b = n-2,\, n-1,\, n\); there are exactly \(3+2+1 = 6\) such points. Translate the situation by \((0,\,-(n-3))\): the map\par
\[
(x,\,y) \;\longmapsto\; (x,\,y-(n-3))
\]
sends these six points bijectively onto the set \(P_3\) (because \(a+b \le n+1\) becomes \(a + (b'+(n-3)) \le n+1 \iff a+b' \le 4\)). Therefore, if we take the three sunny lines that cover \(P_3\) in the translated coordinates and then translate them back, we obtain three sunny lines in the original plane:\par
\[
\begin{aligned}
S_1 &: y = x + (n-3),\\
S_2 &: y = -2x + (n+2),\\
S_3 &: y = -\tfrac12 x + \tfrac{2n-1}{2}.
\end{aligned}
\]
Together with the \(n-3\) horizontal lines we have \(n\) distinct lines, exactly three of which are sunny. Hence \(k = 3\) is attainable for every \(n\ge 3\).\par
\smallskip
\solutionheading{2. Upper bound on points per sunny line}
\textbf{Lemma 1.} For any sunny line \(L\), we have \(|L \cap P_n| \le \left\lfloor\dfrac{n+1}{2}\right\rfloor\).\par
\emph{Proof.} If \(L\) contains no integer point, then obviously \(|L\cap P_n|\le 1\le\lfloor(n+1)/2\rfloor\) (since \(n\ge 3\)). So assume \(L\) contains at least one integer point. Since \(L\) is sunny, it is not vertical, so we can write its equation as \(y = mx + d\).\par
\compactbullet{If \(m\) is irrational, then \(L\) can contain at most one integer point (two integer points would give a rational slope). Hence the bound holds. Thus we may assume \(m\) is rational.}
Write \(m = p/q\) in lowest terms, \(q>0\), \(\gcd(p,q)=1\). Choose an integer point \((x_0,y_0)\) on \(L\). Then all integer points on \(L\) are\par
\[
(x,y) = (x_0 + q t,\; y_0 + p t),\qquad t\in\mathbb Z.
\]
Let the set of \(t\) for which the point lies in \(P_n\) be a consecutive integer interval \(t_{\min},\dots,t_{\max}\); let \(N\) be the number of such \(t\). Then \(t_{\max}-t_{\min} = N-1\).\par
For any such point we have the constraints\par
\[
1 \le x \le n,\qquad 1 \le y \le n,\qquad 2 \le x+y \le n+1.
\]
In particular,\par
\begin{align*}
q\,(N-1) &= |x_{\max}-x_{\min}| \le n-1, \quad\text{(A)}\\
|p|\,(N-1) &= |y_{\max}-y_{\min}| \le n-1, \quad\text{(B)}\\
|p+q|\,(N-1) &= |(x+y)_{\max}-(x+y)_{\min}| \le n-1. \quad\text{(C)}
\end{align*}
Set \(M = \max\{q,\ |p|,\ |p+q|\}\). Then (A)-(C) give\par
\[
N-1 \le \frac{n-1}{M}.
\]
Now, because \(L\) is sunny, we have \(p \neq 0\) (otherwise horizontal) and \(p+q \neq 0\) (otherwise slope \(-1\)). We claim that \(M \ge 2\). Suppose \(M = 1\). Then \(q = 1\) (since \(q>0\)), \(|p| \le 1\), \(|p+q| \le 1\). As \(p \neq 0\), either \(p = 1\) or \(p = -1\).\par
\compactbullet{If \(p = 1\), then \(|p+q| = |1+1| = 2 > 1\), contradiction.}
\compactbullet{If \(p = -1\), then \(p+q = 0\), which is forbidden.}
Hence \(M \ge 2\). Consequently,\par
\[
N-1 \le \frac{n-1}{2}\quad\Longrightarrow\quad N \le \frac{n+1}{2}.
\]
Since \(N\) is an integer, \(N \le \left\lfloor\frac{n+1}{2}\right\rfloor\). \(\square\)\par
\smallskip
\solutionheading{3. No covering by all sunny lines for \(n \ge 4\)}
\textbf{Lemma 2.} For \(n \ge 4\) there is no set of \(n\) distinct sunny lines that covers \(P_n\).\par
\emph{Proof.}\par
\emph{Even \(n\):} Write \(n = 2m\). By Lemma 1 each sunny line contains at most \(m\) points. Therefore the total number of incidences (counting each point as many times as lines through it) is at most \(n \cdot m = 2m^2\). But\par
\[
|P_n| = \frac{n(n+1)}{2} = m(2m+1) = 2m^2 + m > 2m^2 \qquad (m\ge1),
\]
so it is impossible to cover all points.\par
\emph{Odd \(n\):} Write \(n = 2m+1\) with \(m \ge 2\) (i.e. \(n \ge 5\)). Lemma 1 gives each sunny line at most \(m+1\) points. Hence\par
\[
\sum_{i=1}^{n} |L_i \cap P_n| \le n (m+1) = (2m+1)(m+1).
\]
But\par
\[
|P_n| = \frac{(2m+1)(2m+2)}{2} = (2m+1)(m+1).
\]
Thus the inequality must be an equality. In particular, each line contains exactly \(m+1\) points, and the sum of the sizes equals \(|P_n|\). Since every point is covered at least once, equality forces each point to be covered exactly once, and therefore the \(n\) lines are pairwise disjoint.\par
We now classify all sunny lines that can contain exactly \(m+1\) points in \(P_n\).\par
\textbf{Lemma 3.} Let \(n = 2m+1 \ge 3\). If a sunny line \(L\) satisfies \(|L \cap P_n| = m+1\), then \(L\) is one of the three lines\par
\[
L_A: y = x,\qquad L_B: y = -2x + (n+2),\qquad L_C: y = -\tfrac12 x + \tfrac{n+2}{2}.
\]
\emph{Proof.} Write \(L\) in reduced form \(y = \frac{p}{q}x + d\) with \(q>0\), \(\gcd(p,q)=1\), \(p \neq 0\), \(p+q \neq 0\). Let \(N = m+1\).\par
From (A)-(C) we have (with \(N-1 = m\) and \(n-1 = 2m\))\par
\[
q\,m \le 2m,\quad |p|\,m \le 2m,\quad |p+q|\,m \le 2m,
\]
hence\par
\[
q \le 2,\quad |p| \le 2,\quad |p+q| \le 2.
\]
We examine the possibilities.\par
\textbf{Case }\(q = 1\).\textbf{ Then }\(|p| \le 2\) and \(|p+1| \le 2\). Since \(p \neq 0\) and \(p+1 \neq 0\) (i.e. \(p \neq -1\)), the admissible integers are \(p = 1\) and \(p = -2\).\par
\textbf{Case }\(q = 2\).\textbf{ Then }\(\gcd(p,2)=1\), so \(p\) is odd. \(|p| \le 2\) forces \(p = \pm 1\). \(|p+2| \le 2\) eliminates \(p = 1\) because \(|3| > 2\); thus only \(p = -1\) remains.\par
Now we determine the intercept \(d\) so that exactly \(m+1\) integer points of \(P_n\) lie on the line.\par
\smallskip
\textbf{Subcase (1, 1):} \(L: y = x + d\).\par
Let the integer points be \((a+i,\ a+i+d)\) for \(i = 0,1,\dots,m\) (since \(x\) increases by \(q=1\)). The \(x\)-coordinates range from \(a\) to \(a+m\) and must satisfy \(1 \le a+i \le 2m+1\). Hence \(1 \le a \le m+1\). The constraints from the bounds on \(y\) and on \(x+y\) yield:\par
\[
\begin{cases}
a+d \ge 1,\\
a+m+d \le 2m+1,\\
2a+2m+d \le 2m+2.
\end{cases}
\]
The third inequality gives \(d \le 2-2a\); the second gives \(d \le m+1 - a\); the first gives \(d \ge 1-a\).\par
For consistency we need \(1-a \le 2-2a\), i.e. \(a \le 1\). Thus \(a = 1\). Then the bounds become \(0 \le d \le \min\{m,\,0\} = 0\), so \(d = 0\). Hence \(L: y = x\).\par
\smallskip
\textbf{Subcase (1, -2):} \(L: y = -2x + d\).\par
Again the points are \((a+i,\ -2(a+i)+d)\) with \(i=0,\dots,m\) and \(1 \le a \le m+1\). The tightest constraints come from \(y \ge 1\) (at \(i=m\)), \(y \le 2m+1\) (at \(i=0\)), and \(x+y \le 2m+2\) (at \(i=0\)):\par
\[
\begin{aligned}
-2(a+m)+d \ge 1 \;&\Longrightarrow\; d \ge 1 + 2a + 2m,\\
-2a + d \le 2m+1 \;&\Longrightarrow\; d \le 2m+1 + 2a,\\
-a + d \le 2m+2 \;&\Longrightarrow\; d \le 2m+2 + a.
\end{aligned}
\]
The first two give\par
\[
1+2a+2m \le d \le 2m+1+2a.
\]
Hence \(d\) must equal \(2m+1+2a\) (the only integer satisfying both). Substituting into the third inequality yields\par
\[
2m+1+2a \le 2m+2 + a \;\Longrightarrow\; a \le 1.
\]
Since \(a \ge 1\), we have \(a = 1\). Then \(d = 2m+1+2 = 2m+3 = n+2\). So \(L: y = -2x + (n+2)\).\par
\smallskip
\textbf{Subcase (2, -1):} \(L: y = -\tfrac12 x + d\).\par
The \(x\)-coordinates differ by \(q=2\). Let the smallest be \(a\). Then the points are \((a+2i,\ -\tfrac12(a+2i)+d)\) for \(i=0,\dots,m\). To stay within the bounds \(1 \le x \le 2m+1\) we need\par
\[
a \ge 1 \quad\text{and}\quad a+2m \le 2m+1 \;\Longrightarrow\; a \le 1,
\]
so \(a = 1\). Thus the \(x\)-coordinates are \(1,3,5,\dots,2m+1\).\par
For the \(y\)-coordinates to be integers, \(d\) must have the form \(d = d' + \tfrac12\) with \(d' \in \mathbb Z\). Then\par
\[
y_i = -\tfrac12(2i+1) + d = -i - \tfrac12 + d' + \tfrac12 = d' - i.
\]
Now apply the constraints:\par
\compactbullet{\(y_i \ge 1\) for all \(i\): the smallest \(y\) occurs at \(i=m\) (since slope negative), so \(d' - m \ge 1 \;\Rightarrow\; d' \ge m+1\).}
\compactbullet{\(y_i \le 2m+1\) for all \(i\): the largest \(y\) occurs at \(i=0\), giving \(d' \le 2m+1\).}
\compactbullet{\(x_i + y_i = (2i+1) + (d' - i) = d' + i + 1 \le 2m+2\): the most restrictive is at \(i=m\), giving \(d' + m + 1 \le 2m+2 \;\Rightarrow\; d' \le m+1\).}
Thus \(d' \ge m+1\) and \(d' \le m+1\), so \(d' = m+1\). Then \(d = m+1 + \tfrac12 = \tfrac{2m+3}{2} = \tfrac{n+2}{2}\). Hence \(L: y = -\tfrac12 x + \tfrac{n+2}{2}\).\par
\smallskip
We have exhausted all possibilities; therefore any sunny line with exactly \(m+1\) points must be one of the three listed. \(\square\)\par
Returning to the proof of Lemma 2 for odd \(n\ge 5\): we have \(n\) distinct sunny lines, each must contain exactly \(m+1\) points, hence each must be one of \(L_A, L_B, L_C\). But there are only three distinct such lines, while \(n \ge 5\) gives a contradiction. Hence no covering exists for odd \(n\ge 5\) either. (For \(n=3\) the three lines are exactly these three, and that case is covered by the construction.) This completes the proof of Lemma 2. \(\square\)\par
\smallskip
\solutionheading{4. Induction on \(n\)}
We now prove by strong induction that for every \(n\ge 3\) the number \(k\) of sunny lines in any covering of \(P_n\) by \(n\) distinct lines can only be \(0,1,3\).\par
\textbf{Base case }\(n = 3\).\textbf{ Constructions in 1 show that }\(k = 0,1,3\) are attainable. We must show that \(k = 2\) is impossible.\par
Assume three lines cover \(P_3\) with exactly two sunny. Let the sunny lines be \(L_1, L_2\) and the third line \(L_3\) be non-sunny.\par
Lemma 1 implies that each sunny line contains at most two points of \(P_3\). The only pairs of points in \(P_3\) that are collinear with a sunny line (i.e. with slope different from \(0\), undefined, or \(-1\)) are exactly the three pairs:\par
\[
\{(1,1),(2,2)\}\ (\text{slope }1),\quad \{(1,3),(2,1)\}\ (\text{slope }-2),\quad \{(1,2),(3,1)\}\ (\text{slope }-1/2).
\]
Hence any sunny line that contains two points must be one of the three lines\par
\[
\ell_1: y=x,\quad \ell_2: y=-2x+5,\quad \ell_3: y=-\tfrac12 x+\tfrac52.
\]
Now consider the possibilities.\par
\emph{If both \(L_1\) and \(L_2\) contain two points and are disjoint.} Then they must be two of the three lines \(\ell_1,\ell_2,\ell_3\). Those three lines are pairwise disjoint, so they cover four points. The two remaining points lie on the third \(\ell\), which is sunny. The third line \(L_3\) is non-sunny, so it cannot cover both of those points (the unique line through them is sunny). Contradiction.\par
\emph{If at least one of \(L_1, L_2\) contains at most one point.} Then the two sunny lines together cover at most three points. Consequently, the non-sunny line \(L_3\) must cover at least three points. The only lines that contain three points of \(P_3\) are the three "boundary" lines:\par
\[
y=1,\quad x=1,\quad x+y=4.
\]
(One verifies that any other line contains at most two points.) So \(L_3\) is one of these three.\par
\compactbullet{\textbf{}\(L_3 = y=1\)\textbf{: covers }\((1,1),(2,1),(3,1)\). The uncovered points are \((1,2),(1,3),(2,2)\). Each of these three points lies on a different one of \(\ell_1,\ell_2,\ell_3\) (namely \((2,2)\in\ell_1\), \((1,3)\in\ell_2\), \((1,2)\in\ell_3\)). Two sunny lines can cover at most two of them (each \(\ell_i\) contains exactly one). Hence impossible.}
\compactbullet{\textbf{}\(L_3 = x=1\)\textbf{: covers }\((1,1),(1,2),(1,3)\). Uncovered: \((2,1),(2,2),(3,1)\). Again these are one per \(\ell_i\); two sunny lines cannot cover all three. Contradiction.}
\compactbullet{\textbf{}\(L_3 = x+y=4\)\textbf{: covers }\((1,3),(2,2),(3,1)\). Uncovered: \((1,1),(1,2),(2,1)\), again one per \(\ell_i\). Contradiction.}
Thus \(k = 2\) is impossible, so the only possible values for \(n=3\) are \(0,1,3\).\par
\textbf{Inductive step.} Assume the statement holds for all \(m\) with \(3 \le m < n\), where \(n \ge 4\). Consider any covering of \(P_n\) by \(n\) distinct lines.\par
Define three special lines:\par
\[
R:\ y=1,\qquad C:\ x=1,\qquad D:\ x+y = n+1.
\]
Note that \(R, C, D\) are not sunny (horizontal, vertical, slope \(-1\) respectively).\par
\textbf{Claim.} At least one of \(R, C, D\) belongs to the covering.\par
\emph{Proof of claim.} Suppose none of them is present. Then:\par
\compactbullet{The line \(y=1\) contains \(n\) distinct points \((a,1)\). To cover a point \((a,1)\), the line must intersect \(y=1\) at that point. A horizontal line different from \(y=1\) does not intersect \(y=1\). Hence every line that covers a point on \(y=1\) must be non-horizontal. Moreover, a non-horizontal line meets \(y=1\) in at most one point. Therefore we need at least \(n\) non-horizontal lines; since there are exactly \(n\) lines, \emph{all} lines are non-horizontal, and each covers exactly one point on \(y=1\).}
\compactbullet{The line \(x=1\) contains \(n\) points \((1,b)\). Since \(C\) is absent, a vertical line different from \(x=1\) does not contain any point with \(x=1\). Hence every line covering a point on \(x=1\) must be non-vertical. Since all lines are already non-horizontal, they must all be non-vertical to cover those points (a vertical line would miss \(x=1\) unless it is \(x=1\) itself). Thus all lines are non-vertical.}
\compactbullet{The line \(x+y=n+1\) contains \(n\) points. Since \(D\) is absent, a line with slope \(-1\) other than \(D\) is parallel to it and does not intersect it. Therefore every line covering a point on \(D\) must have slope different from \(-1\). Hence all lines have slope \(\neq -1\).}
Therefore every line is non-horizontal, non-vertical, and slope \(\neq -1\); i.e., every line is sunny. But Lemma 2 states that for \(n\ge 4\) there is no covering of \(P_n\) by \(n\) sunny lines. Contradiction. Hence the claim holds. \(\square\)\par
Now we examine which of \(R, C, D\) is present. (If more than one is present, we may choose any one; the argument works in each case.)\par
\emph{If \(R\) is present:} Remove \(R\) from the set. The remaining \(n-1\) lines still cover all points of \(P_n\) with \(y\ge 2\) (because \(R\) only covered points with \(y=1\)). Apply the translation \(\phi_R: (x,y) \mapsto (x, y-1)\). This map sends the set \(\{(x,y)\in P_n \mid y\ge 2\}\) bijectively onto \(P_{n-1}\). Moreover, \(\phi_R\) sends each line to a line (translations preserve lines). The resulting family consists of \(n-1\) distinct lines (translation preserves distinctness) and covers \(P_{n-1}\) (each point of \(P_{n-1}\) corresponds to a point with \(y\ge 2\) that was covered). Since \(R\) is not sunny, the number of sunny lines among the \(n-1\) lines remains \(k\).\par
Thus we have a covering of \(P_{n-1}\) by \(n-1\) distinct lines with exactly \(k\) sunny lines. By the induction hypothesis (applied to \(m=n-1\)), we must have \(k \in \{0,1,3\}\).\par
\emph{If \(C\) is present:} Remove \(C\) and use \(\phi_C: (x,y) \mapsto (x-1, y)\). It maps \(\{(x,y)\in P_n \mid x\ge 2\}\) onto \(P_{n-1}\). \(C\) is non-sunny, so \(k\) unchanged. Induction gives \(k \in \{0,1,3\}\).\par
\emph{If \(D\) is present:} Remove \(D\). The remaining \(n-1\) lines cover all points of \(P_n\) that are not on \(D\). Observe that \(P_{n-1} \subseteq P_n\) and \(P_{n-1} \cap D = \emptyset\) (since points in \(P_{n-1}\) satisfy \(a+b\le n\), while \(D\) requires \(a+b = n+1\). Hence they cover \(P_{n-1}\). \(D\) is not sunny, so \(k\) stays the same. Induction yields \(k \in \{0,1,3\}\).\par
In every possible case we conclude \(k \in \{0,1,3\}\), which completes the inductive step.\par
By strong induction, the statement holds for all integers \(n \ge 3\).\par
\smallskip
\solutionheading{5. Conclusion}
We have shown that for every \(n \ge 3\) the only possible numbers of sunny lines in a covering of \(P_n\) by \(n\) distinct lines are \(0\), \(1\), and \(3\). Explicit constructions demonstrate that each of these values is indeed attainable.\par
\[
\boxed{k \in \{0,\,1,\,3\}\ \text{ for all } n\ge 3.}
\]
\par
